\section{Related Work}
\label{sec:related}

\paragraph{Network facility location.}
The \mbox{p-median} problem seeks to select $p$ service sites that minimize
total demand-weighted assignment distance~\cite{hakimi1964optimum}.
Its network variant, in which distances are shortest paths on a graph rather
than Euclidean measures, is the natural model when travel effort matters.
Maranzana's early observation that Euclidean distances understate network
costs in irregular urban street layouts~\cite{maranzana1964siting} has been
confirmed repeatedly in empirical studies; Boeing's OSMnx library has made
pedestrian-network extraction operationally straightforward~\cite{boeing2017osmnx}.
Solution methods for medium-to-large instances range from Lagrangian relaxation
and column generation to commercial MILP
solvers~\cite{avella2007lagrangian,lorena2003constructive}.
For instances with thousands of candidate sites and demand nodes,
decomposition approaches that reduce the effective problem size have
practical value, as long as the decomposition does not introduce large
coverage gaps relative to the global optimum.

\paragraph{Clustering-based territorial decomposition.}
Partitioning a study area before solving a subproblem independently in each
territory is a well-known heuristic for scaling facility location to larger
instances.
The quality of the decomposition depends on whether the partition respects
the connectivity of the underlying network.
Centroid-based methods such as \mbox{k-means} impose Euclidean distances on
both the partition and the assignment, which introduces inconsistency when
actual travel distances are non-Euclidean.
Spectral clustering~\cite{shi2000normalized,ng2002spectral,vonluxburg2007tutorial}
avoids this by using the normalized graph Laplacian of an affinity matrix
defined on road-network distances; the resulting partition follows the
connectivity structure of the street network rather than Euclidean geometry.
The eigengap heuristic provides a data-driven criterion for selecting the
number of clusters, and a silhouette-score secondary check supplies additional
validation~\cite{theodoridis2009clustering}.
A post-processing merge rule that absorbs small territories into their nearest
major cluster bridges the gap between the mathematically selected partition
and the operationally deployable territory count.

\paragraph{Socioeconomic demand weighting.}
Demand models in urban facility location typically scale by population count,
which disregards heterogeneity in purchasing power, infrastructure quality,
and transport access.
MCA provides a principled method for compressing multiple categorical census
variables into orthogonal latent dimensions~\cite{greenacre1988correspondence,%
greenacre2006multiple,greenacre2017correspondence}.
Because census data are inherently categorical or categorized, MCA is preferable
to PCA: it operates on the indicator matrix of dummy-coded categories rather
than on a covariance matrix of continuous variables.
Greenacre's corrected-inertia formula removes the inflation that raw MCA
eigenvalues exhibit when the number of indicator columns is large, making
dimension importance comparable across models~\cite{greenacre1991interpreting}.
The normalized first MCA dimension serves as a socio-urban quality gradient
that upweights blocks with better infrastructure, education, and housing
services, directing new openings toward areas where economic activity
is most likely to be sustained.

\paragraph{Gap in existing work.}
Prior applications of \mbox{p-median} models to convenience-store or retail
location typically fix parameters by convention, use Euclidean distances as
a proxy, compare against no baseline, and validate in a single geographic
context.
This paper addresses all four gaps: parameters are calibrated from data,
distances are pedestrian-network distances in both solving and evaluation,
four controlled baselines are compared on equal terms, and the framework
is validated in three urban settings with qualitatively different
supply-demand structures.
